% !TeX spellcheck = en_US
\documentclass[french]{yLectureNote}

\title{Électromagnétisme 1 PS}
\subtitle{Physique}
\author{Paulhenry Saux}
\date{\today}
\yLanguage{Français}
%lionels.calmels@cemes.fr
\professor{F.Pettinari}
\usepackage{graphicx}%----pour mettre des images
\usepackage[utf8]{inputenc}%---encodage
\usepackage{geometry}%---pour modifier les tailles et mettre a4paper
%\usepackage{awesomebox}%---pour les boites d'exercices, de pbq et de croquis ---d\'esactiv\'e pour les TP de PC
\usepackage{tikz}%---pour deiffner + d\'ependance de chemfig
% \usepackage{tabularx}%---pour dimensionner automatiquement les tableaux avec variable X
\usepackage{awesomebox}%---Pour les boites info, danger et autres
\usepackage{menukeys}%---Pour deiffner les touches de Calculatrice
\usepackage{fancyhdr}%---pour les en-t\^ete personnalis\'ees
\usepackage{blindtext}%---pour les liens
\usepackage{hyperref}%---pour les liens (\`a mettre en dernier)
\usepackage{caption}%---pour la francisation de la l\'egende table vers Tableau
\usepackage{pifont}
\usepackage{array}%---pour les tableaux
\usepackage{lipsum}
\usepackage{yFlatTable}
\usepackage{multicol}
\newcommand{\Lim}[1]{\lim\limits_{\substack{#1}}\:}
\renewcommand{\vec}{\overrightarrow}
\newcommand{\N}[0]{\mathbb{N}}
\newcommand{\dd}{\mathrm{d}}
\newcommand{\norm}[1]{||\vec{#1}||}
\newcommand{\fo}{\psi(\vec{r},t)}
\newcommand{\foe}{\psi(\vec{r},t)\*}
\newcommand{\HH}{\hat{H}}
\newcommand{\hb}{\hbar}
\newcommand{\lap}{\nabla^2}
\newcommand{\lapcc}{\frac{\partial^2 }{\partial x^2}+\frac{\partial^2 }{\partial y^2}+\frac{\partial^2 }{\partial z^2}}
\newcommand{\E}{\vec{E}(M)}
\newcommand{\V}{V(M)}
\newcommand{\er}{\vec{e_{\rho}}}
\newcommand{\ep}{\vec{e_{\varphi}}}
\newcommand{\pip}{\pi^+}
\newcommand{\pim}{\pi^-}
\newcommand{\mv}{\mathcal{V}}
\begin{document}
\setcounter{chapter}{4}
\chapter{Calcul de champ avec Gauss}
Cette méthode ne s'applique que dans le cas où le système présente beaucoup d'éléments de symétrie\marginCheck{\begin{itemize}
 \item C\^able rectiligne infiniment long tel que \(\rho_c(M) = \rho_c(\rho)\)
 \item Distributions de charge à symétrie sphérique
 \item Plaques infinies planes \(\perp (Oz)\) telles que \(\rho_c(M) = \rho_c(z)\).
\end{itemize}} :

\section{Notion de surface fermée orienté}
\begin{definition}[Surface fermée]
C'est une surface qui délimite entièrement un volume \(\mv\) fini.
\end{definition}
C'est le cas des sphères, boite de conserve, boite cubique. Des surfaces ouvertes sont des carrés, paroi latérale d'un cylindre, disque.

\begin{definition}[Surface fermée orientée]
Une surface fermée orientée  est une surface fermée en tout de laquelle on associe un vecteur unitaire \(\vec{n}(M)\) perpendiculaire en ce point à la surface et orienté vers l'extérieur de la surface.
\end{definition}
\section{Flux de champ}
En un point M, on définit la surface élémentaire autour de ce point et le vecteur normal à la surface. On définit également le champ électrique crée par la distribution de charge.

\begin{definition}
Le flux de E à travers S est le scalaire \[\phi = \iint \E \cdot \vec{n}\dd S\]
\end{definition}
% \section{Calcul de flux}
% \subsection{Généralisation}
% Si E est le champ crée par une charge ponctuelle \(q\) alors \(iint \E \cdot \vec{n}\dd S = \frac{q}{\epsilon_0}\) ou 0 si elle ne la contient pas.
\section{Théorème de Gauss}
On considère une distribution de charges et une surface fermée S. Une partie est à l'intérieur de S et l'autre à l'extérieur.

On note \(\E\) le champ électrique crée au point M par toutes les charges de la distribution, \(E_{int}, E_{ext}\) respectivement celui crée par les charges intérieures et extérieures.

% On sait que \[\iint \E\cdot \vec{n}\dd S = \iint E_{int}\cdot \vec{n}\dd S = \sum q_{int} \cdot \frac{1}{\epsilon_0}\]
\begin{theorem}[Théorème de Gauss]
Le flux de E à travers une surface fermée orientée est égal à la somme des charges intérieures à cette surface fermée divisée par \(\epsilon_0\) :
\[\iint \E\cdot \vec{n}\dd S = \frac{\sum q_{int}}{\epsilon_0}\]
\end{theorem}
\section{Méthode de calcul}
\begin{enumerate}
 \item Faire un choix judicieux, d'axes, de système de coordonnées et de base.
 \item Considérer les symétries et invariances pour savoir quelles sont les composantes non nulle et leur dépendance.
 \item Faire un choix judicieux de surface fermée S (Surface de Gauss) servant à appliquer le théorème de Gauss
 \begin{itemize}
  \item Elle doit \^etre telle que \(\E\cdot \vec{n}\) est uniforme sur S
 \end{itemize}
\item Expliciter le calcul du Flux en fonction de la composante non nulle du champ électrique
\item Calculer la somme des charges intérieures à S
\item Appliquer le théorème de Gauss et en déduire la composante inconnue du champ
\end{enumerate}
\section{Exemples}
\subsection{Charge ponctuelle}
La distribution de charge est àn symétrie sphérique autour de la charge ponctuelle placée en O. On utilise les coordonnées sphériques.

On a \(\E = E_r(r)\er\)

On choisit maintenant S : une sphère de centre O pour obtenir que \[\E\cdot \vec{n} = E_r(r)\]
On fait maintenant le calcul explicite: \[\iint \E\cdot \vec{n}\dd S = E_r\cdot r^2\int_{\theta=0}^{\pi}\sin(\theta)\dd \theta\int_{\varphi=0}^{2\pi} = 4\pi r^2E_r\]
On sait que \(\sum q_{int} = q\)

Donc par le théorème de Gauss :
\[E_r  = \frac{q}{4\pi\epsilon_0r^2}\er\]
\subsection{Fil rectiligne infini}
On confond l'axe Oz avec un fil chargé et on utilise la base cylindrique locale associée à M. De plus, par les symétries et invariances, \(\E = E_{\rho}\vec{e_{\rho}}\).
\explanation{1}{Si on prend un point quelconque sur \(S_1\) et \(S_2\), le produit scalaire est bien uniforme par morceaux sur les 3 faces. En plus, le flux sur les faces haut et bas s'annule.}

Choix de S : Cylindre fermé d'axe Oz de rayon \(\rho\)\explain{1}{right}{0}{0.5}{}.
\begin{flalign*}
\phi &= \iint \E\cdot\vec{n}\dd S\\
&= \iint_{S1}\E\cdot \vec{n}\dd S + \iint_{S2}\E\cdot \vec{n}\dd S+ \iint_{Slat}\E\cdot \vec{n}\dd S\\
&= 0 + 0 + \iint_{Slat} E_{\rho}(\rho)\rho\dd \varphi\dd z\\
&= \int_{\varphi=0}^{2\pi}\int_{z=z_1}^{z_2}E_{\rho}(\rho)\rho\dd \varphi\dd z\\
&=E_{\rho}(\rho)\rho \int_{\varphi=0}^{2\pi}\dd \varphi \int_{z=z_1}^{z_2}\dd z\\
&= E_{\rho}(\rho)\rho\cdot2\pi (z_2-z_1)
\end{flalign*}
On calcule maintenant \(\sum Q_{int}\) : \(Q_{int} = (z_2-z_1)\lambda\)

On écrit maintenant le théorème de Gauss :
\begin{flalign*}
2\pi\rho(z_2-z_1)E_{\rho} &= \frac{1}{\epsilon_0}(z_2-z_1)\lambda\\
E_{\rho} &= \frac{\lambda}{2\pi\epsilon_0 \rho}
\end{flalign*}
 \end{document}
